% authority: science-draft
% status: proposal-only
% route: ontology-law-research-packet
% trigger_classification: derivation_critical_missing_source_law
% adoption_status: blocked_adoption_open_continuation
\documentclass[11pt]{article}
\usepackage[margin=1in]{geometry}
\usepackage{amsmath,amssymb,amsthm}
\usepackage{booktabs}

\newtheorem{theorem}{Theorem}
\newtheorem{proposition}{Proposition}
\newtheorem{remark}{Remark}
\newcommand{\Ftwo}{\mathbb F_2}

\title{Refuter Stress of
EqSrcOrderedMotionRootedPartitionLaw\(_{\rm src}^{\rm cand,v1}\):\\
Root-Side Ambiguity, Exhaustive Transition Nonuniqueness, and Variation Orbit}
\author{Refuter control packet RT-20260718-045}
\date{2026-07-19}

\begin{document}
\maketitle

\section{Control result}

This packet stresses the unchanged exact candidate at
\texttt{24992d4b41d64bba860f5cd61d505d6b1ecaad3917e9195e0cbd3d897d955aef}.
Its conditional algebraic core survives: after a finite transition record,
projection, singleton root, morphism category, and variation family are
supplied and admitted, path parity, the ordered partition, the line
\(B_U\), and the kernel-pair relation \(R_U\) are well-defined.

The exact Refuter classification is
\[
  \texttt{scoped\_obstruction}.
\]
The failed premise is not the kernel-pair theorem. It is the premise that
current ontology or the unchanged candidate derives source data selecting a
unique relation. The candidate remains \texttt{proposal-only},
\texttt{new\_ontology\_primitive\_candidate}, and
\texttt{blocked\_adoption\_open\_continuation}.

\section{Root-side orientation lemma}

Let \(G=(X,D)\) be a finite directed graph with a root \(r\) from which every
vertex is reachable. Suppose all directed root-to-vertex paths have a fixed
parity. Then every directed edge crosses the parity blocks \(P_r,Q_r\).
The underlying reachable graph is therefore bipartite.

\begin{theorem}[Root-side orientation]
For a fixed reachable bipartite transition graph, any two structurally
admissible roots in the same bipartition side induce the same ordered blocks.
Roots in opposite sides exchange the blocks. Consequently the graph supplies
at most two candidate boundary lines,
\[
 B=\Ftwo\mathbf 1_Q,\qquad B^\vee=\Ftwo\mathbf 1_P.
\]
When both blocks are nonempty, these lines and their kernel-pair relations are
distinct.
\end{theorem}

\begin{proof}
A connected bipartite graph has a unique two-coloring up to exchange.
Root parity selects the color containing the root as even. Moving the root
within that color preserves the orientation; moving it to the other color
exchanges \(P\) and \(Q\). Since
\(\mathbf 1_P\neq\mathbf 1_Q\), their nonzero one-dimensional spans in
\(\operatorname{span}_{\Ftwo}\{\mathbf 1_P,\mathbf 1_Q\}\) are distinct.
\end{proof}

\begin{remark}
A loaded singleton root chooses one orientation and makes the conditional
candidate unambiguous. The theorem shows why an independent source-derived
root law is a substantive missing primitive rather than notation.
\end{remark}

\section{Exhaustive four-token transition census}

The validator enumerates all \(2^{12}=4096\) loop-free directed graphs on four
labeled tokens and all \(4\) root choices, for \(16384\) graph--root
configurations. A configuration is counted as structurally admitted exactly
when every token is root reachable, all root paths have one parity, and both
parity blocks are nonempty. Transition-certificate provenance is granted as a
proposal assertion; the census does not manufacture physical provenance.

\begin{center}
\begin{tabular}{lr}
\toprule
Exact quantity & Count \\
\midrule
structurally admitted rooted records & 1088 \\
transition graphs admitting at least one root & 557 \\
graphs inducing one relation & 292 \\
graphs whose roots induce two relations & 265 \\
\bottomrule
\end{tabular}
\end{center}

The admitted-graph root-count distribution is \(292,108,48,109\) for one,
two, three, and four roots. The edge-count distribution is
\[
3{:}64,\quad 4{:}186,\quad 5{:}192,\quad 6{:}88,\quad
7{:}24,\quad 8{:}3.
\]

For odd-block sizes \(1,2,3\), the total rooted-record counts are
\(96,960,32\). Each fixed partition of those sizes has respectively
\(24,160,8\) rooted transition realizations. Hence relation selection does not
select unique transition provenance even after the root-oriented partition is
fixed.

\section{Minimal and symmetric countermodels}

The minimal two-token directed cycle has the swap automorphism and no
automorphism-invariant singleton root. Loading root \(0\) gives
\(Q=\{1\}\); loading root \(1\) gives \(Q=\{0\}\). The two boundary lines and
kernel relations are distinct.

The strengthened directed four-cycle
\[
  0\longrightarrow1\longrightarrow2\longrightarrow3\longrightarrow0
\]
has four structurally admissible roots. Roots \(0,2\) give
\(Q=\{1,3\}\), while roots \(1,3\) give \(Q=\{0,2\}\). Its order-four
automorphism group acts transitively on roots. No singleton root is invariant;
only the order-two rotation subgroup stabilizes either selected relation.

These controls do not contradict the candidate's explicit fail-closed branch
for a nonnatural root. They show that the branch is reached unless a new
source-side root-provenance primitive breaks the symmetry. Selecting the first
observed token would do so by observer truncation, which the candidate
correctly forbids.

\section{Complete four-token relation-variation orbit}

There are fourteen nonempty proper odd blocks \(Q\subset\{0,1,2,3\}\).
Every distinct ordered pair \(Q\to Q'\) changes the boundary line. The
\(14\cdot13=182\) changes split by Hamming distance as
\[
\begin{array}{c|rrrr}
|Q\triangle Q'| & 1 & 2 & 3 & 4\\
\hline
\text{ordered changes} & 48 & 72 & 48 & 14.
\end{array}
\]
Thus the earlier \(48\) single-token transfers are only the first layer of
the changing orbit. The unchanged candidate admits none of the \(182\) cases
as relation-preserving variations.

The permutation control independently recovers \(14\cdot24=336\) cases:
\(72\) stabilize the selected boundary and \(264\) change it. The stabilizer
orders are \(6,4,6\) for odd-block sizes \(1,2,3\).

\section{Failure-branch classification}

\begin{center}
\begin{tabular}{lll}
\toprule
Branch & Result & Exact scope \\
\midrule
algebraic collapse & survives & two-class kernel pair on admitted record \\
selector collapse & obstructed & unrestricted records realize all relations \\
nonuniqueness & obstructed & 265 graphs support two root-induced relations \\
inverse defect & none found & admitted source isomorphisms only \\
cocycle defect & not applicable & no cocycle theorem is claimed \\
finite variation & fragile & all 182 relation changes excluded \\
physical covariance & unproved & preserving category is loaded \\
\bottomrule
\end{tabular}
\end{center}

The event projection \(\pi=\mathrm{id}\) suffices for every finite control, so
projection declaration does not reduce the relation family. Likewise,
certificate compatibility with schematic \(\Phi_\lambda\) is asserted rather
than derived. Declaration before \(B_U\) is not a provenance certificate.

\section{Obstruction, freeze, and Distance to GR}

The obstruction is local to exact candidate v1 under the current-ontology
source package. Current ontology does not derive a transition, projection,
root, physical morphism, or variation law that selects one candidate relation.
A future conservative source extension may supply such a law; no theorem here
proves that impossible.

General EqSrc remains undischarged. Source ontology, RetainH, GenH,
\(M_{\rm src}\), \(g_{\rm eff}\), matter coupling, Einstein equations,
benchmark promotion, and Gate Chair status are unchanged. The Distance-to-GR
and metric-use ledgers are unchanged.

Exact candidate-v1 construction, fresh audit, and bounded Refuter stress are
complete. The cycle is locally frozen: repeating, repairing, or promoting the
unchanged candidate is barred. This local freeze is not canonical rejection,
global theory rejection, or future source-extension impossibility.

\section{Required next route}

The next packet is one bounded
\texttt{theoretical-continuation-selector@0.1.0}
\texttt{ontology-law-research-packet}. It must select a genuinely
source-derived transition/root provenance law, a source-side irrelevance
theorem, or a distinct scoped no-go question with new mathematical payload.
It may not repeat, repair, adopt, or promote exact candidate v1.

\end{document}
